\section{概要}

Font to Binary Converter は、ドットフォントで描画された文字を
二次元バイナリ配列（0/1）または各種エンコード形式に変換し、
C/C++/Python/JavaScript 等のソースコードに直接埋め込める形式で
出力する Web ツールである。

本ツールは以下の用途を主な想定としている。

\begin{itemize}[leftmargin=*,itemsep=2pt]
  \item マイコン（Arduino, ESP32, RZ/A1 等）の 0.96'' OLED、LCD 等
        小型ディスプレイ上へのフォント描画用ビットマップデータ生成。
  \item LED マトリクスパネル（8×8, 16×16, 32×32）の文字表示用データ。
  \item ゲーム・デモ用の固定幅ドットフォント素材生成。
  \item ピクセルアート作成時の下書きテンプレート。
\end{itemize}

\subsection{主な特徴}

\begin{itemize}[leftmargin=*,itemsep=2pt]
  \item \textbf{ブラウザ単体で動作}: サーバー不要。GitHub Pages で静的配信。
  \item \textbf{2モードUI}:「かんたん」は3ステップ(文字/見た目/用途)で
        初心者でも迷わず変換可能。「詳細」は全設定を直接編集可能。
        モード選択は LocalStorage に自動保存。
  \item \textbf{用途別プリセット}: Arduino/C言語, C言語2D, Python リスト,
        MicroPython bytes, JSON, 記号アートの6種を1クリック切替。
  \item \textbf{フォント同梱}: \texttt{DotGothic16} (16×16) と
        \texttt{Misaki Gothic 2nd} (8×8) を同梱。
  \item \textbf{ユーザーフォント対応}: TTF/OTF/WOFF をその場でアップロード可能。
  \item \textbf{多言語出力}: C, C++, Arduino(PROGMEM), Python, JavaScript,
        TypeScript, Rust, Go, JSON, Markdown, プレーンテキスト, 記号アート。
  \item \textbf{詳細な出力設定}: 2D/3D/フラット/ビットパック構造、
        10/16/2進表記、MSB/LSB、行/列優先、0/1反転など。
  \item \textbf{ピクセルエディタ}: 生成結果を1ドット単位で手動編集可能。
  \item \textbf{設定の自動保存と共有}: LocalStorage へ自動保存、
        URL クエリパラメータによる共有リンク生成。
\end{itemize}

\subsection{動作環境}

\begin{tabularx}{\textwidth}{l X}
\toprule
\textbf{項目} & \textbf{要件} \\
\midrule
ブラウザ & Chrome / Edge / Firefox / Safari の最新版 \\
JavaScript & 有効 \\
Canvas API & \texttt{CanvasRenderingContext2D}, \texttt{getImageData} 対応 \\
FontFace API & カスタムフォント読込に使用（主要ブラウザで対応済み）\\
LocalStorage & 設定の永続化に使用（オプトアウト可）\\
\bottomrule
\end{tabularx}

\newpage
